\documentclass[12pt]{article}
\usepackage[utf8]{inputenc}
\usepackage[T1]{fontenc}
\usepackage{hyperref}
\title{Sasak Language Resources for K-12 Curriculum Development}
\author{Compiled by Claude Sonnet 5.5 (claude-sonnet-5-5) for Andr\'as Kornai}
\date{}
\begin{document}
\maketitle

\section{Introduction}
Sasak (ISO 639-3 \texttt{sas}; Glottocode sasa1249) is an Austronesian language of Lombok, Indonesia, with roughly two million speakers. It has four or five major dialects and a system of speech levels. It was traditionally written in a script close to Balinese, and the Latin alphabet is used today. Indonesian is the language of education. Search results show local-content (\emph{muatan lokal}) instruction in Sasak does exist in places, contrary to the reference description.

\section{Primary Educational Resources}
\begin{itemize}
\item Peraturan Bupati Lombok Tengah (Central Lombok regent's regulation) of 2024 on the local-content curriculum for Sasak as a regional language in basic education: \url{https://jnnqvmfhkqwfdlzevfao.supabase.co/storage/v1/object/public/regulation-pdfs/perbup-kab-lombok-tengah-no-76-tahun-2024.pdf}
\item Report that Balai Bahasa NTB supported the launch of the local-content book \emph{Berajah Base Sasak} for primary-school teachers in Central Lombok: \url{https://rri.co.id/en/iptek/2527933/balai-bahasa-ntb-dukung-peluncuran-buku-muatan-lokal-bahasa-sasak-di-lombok-tengah}
\end{itemize}

\section{Linguistic and Reference Materials}
\begin{itemize}
\item ``Standarisasi Tatabahasa Bahasa Sasak'' (Adabiyyat, UIN Sunan Kalijaga), discussing a standard grammar; it cites \emph{Struktur Bahasa Sasak Umum} (1995) and \emph{Standarisasi Ejaan dan Tata Bahasa Sasak} (2005): \url{https://ejournal.uin-suka.ac.id/adab/Adabiyyat/en/article/view/570}
\item ``Sikap dan pemertahanan bahasa Sasak pada masyarakat Lombok'' (language attitudes and maintenance): \url{https://ejurnal.unim.ac.id/index.php/matapena/article/view/2213}
\item Sasak language overview (Wikipedia): \url{https://en.wikipedia.org/wiki/Sasak_language}
\item Sasak script overview: \url{https://en.wikipedia.org/wiki/Sasak_script}
\item Austronesian's Traces in Sasak: Historical Linguistics Study (Arbitrer): \url{https://arbitrer.fib.unand.ac.id/index.php/arbitrer/article/view/823}
\item Telkom University thesis on an Android Sasak-Indonesian dictionary app (not a published dictionary): \url{https://repositori.telkomuniversity.ac.id/pustaka/files/121197/abstract/aplikasi-kamus-bahasa-sasak-berbasis-android.pdf}
\end{itemize}

\section{Religious and Community Resources}
\begin{itemize}
\item Scripture Earth page for Sasak (Today's Sasak Version New Testament, 2007, and other portions): \url{https://scriptureearth.org/00eng.php?iso=sas}
\end{itemize}

\section{Digital and Audio Resources}
\begin{itemize}
\item Global Recordings Network, Sasak audio: \url{https://globalrecordings.net/en/language/sas}
\item Omniglot page on Sasak writing: \url{https://omniglot.com/writing/sasak.htm}
\end{itemize}

\section{Curriculum Development Context}
Sasak is widely spoken, and local-content teaching is being formalised in Central Lombok. The regulation and the new teachers' book suggest a starting point for primary grades. Dialect variation and speech levels need decisions about which standard to use. I did not find a published Sasak-Indonesian dictionary or a full grammar, so these would need more targeted searching, for instance in the Balai Bahasa NTB catalogue.

\end{document}